\documentclass{article}
\usepackage{amsmath, amssymb, amsthm}
\usepackage{hyperref}
\usepackage{geometry}
\geometry{margin=1in}

\title{Erd\H{o}s Problem \#757}
\date{Last updated: 2025-08-31}
\author{Source: erdosproblems.com}

\begin{document}
\maketitle

\noindent\textbf{Status:} OPEN \quad \textbf{No prize}

\noindent\textbf{Tags:} geometry, distances, sidon sets

\medskip
\noindent\textbf{Problem Statement:}

Let $A\subset \mathbb{R}$ be a set of size $n$ such that every subset $B\subseteq A$ with $\lvert B\rvert =4$ has $\lvert B-B\rvert\geq 11$. Find the best constant $c>0$ such that $A$ must always contain a Sidon set of size $\geq cn$.

\bigskip
\noindent\textbf{Remarks:}

For comparison, note that if $B$ were a Sidon set then $\lvert B-B\rvert=13$, so this condition is saying that at most one difference is 'missing' from $B-B$. Equivalently, one can view $A$ as a set such that every four points determine at least five distinct distances, and ask for a subset with all distances distinct.

Without loss of generality, one can assume $A\subset \mathbb{N}$. 

Erd\H{o}s and S\'{o}s proved that $c\geq 1/2$. Gy\'{a}rf\'{a}s and Lehel \cite{GyLe95} proved\[\frac{1}{2}+\frac{1}{141\cdot 76}\leq c\leq\frac{3}{5}.\](The example proving the upper bound is the set of the first $n$ Fibonacci numbers.)
Ma and Tang \cite{MaTa26} have improved these bounds to\[\frac{9}{17}\leq c\leq \frac{4}{7}.\]

\bigskip
\noindent\textbf{References:}

[GyLe95] Gy\'{a}rf\'{a}s, Andr\'{a}s and Lehel, Jen\H{o}, Linear sets with five distinct differences among any four
elements. J. Combin. Theory Ser. B (1995), 108-118.
 
 [MaTa26] J. Ma and Q. Tang, Largest Sidon subsets in weak Sidon sets. arXiv:2602.23282 (2026).

\bigskip
\noindent\small{Source: \url{https://www.erdosproblems.com/757}}
\end{document}
